\bta{2018}

\begin{enumerate}
	%\renewcommand{\labelenumi}{\arabic{enumi}.}
	% A(\Alph) a(\alph) I(\Roman) i(\roman) 1(\arabic)
	%设定全局标号series=example	%引用全局变量resume=example
	%[topsep=-0.3em,parsep=-0.3em,itemsep=-0.3em,partopsep=-0.3em]
	%可使用leftmargin调整列表环境左边的空白长度 [leftmargin=0em]
	\item
	质量为$ m $的粒子在一维无限深方势阱 $(0<\mathrm{x}<\mathrm{a})$ 中运动， 设初始波函数 $\psi(x, 0)=A\left[\phi_{1}(x)+\phi_{2}(x)\right]$, 其中 $\phi_{1}(x), \phi_{2}(x)$ 分别表示归一化的基态和第一激发态波函数。
	\begin{enumerate}
		%\renewcommand{\labelenumi}{\arabic{enumi}.}
		% A(\Alph) a(\alph) I(\Roman) i(\roman) 1(\arabic)
		%设定全局标号series=example	%引用全局变量resume=example
		%[topsep=-0.3em,parsep=-0.3em,itemsep=-0.3em,partopsep=-0.3em]
		%可使用leftmargin调整列表环境左边的空白长度 [leftmargin=0em]
		\item
		写出 $\phi_{1}(x), \phi_{2}(x)$ 和归一化因子$ A $。
		\item 
		求 $\psi(x, t)$。
		\item 
		计算$ t=0 $时粒子坐标算符x和动量算符p的平均值。
		\item 
		计算$ t>0 $时粒子能量的可取值，对应概率和平均值。
		
		
		
	\end{enumerate}
	
		\banswer{
			\begin{enumerate}
				\item
				\[
					\phi_1(x)=\sqrt{\frac2a}\sin\frac{\pi x}{a},\qquad
					\phi_2(x)=\sqrt{\frac2a}\sin\frac{2\pi x}{a}.
				\]
				由正交归一性，$1=|A|^2(1+1)$，可取
				$\boxed{A=1/\sqrt2}$。

				\item 令
				$E_n=n^2\pi^2\hbar^2/(2ma^2)$，则
				\[
					\boxed{\psi(x,t)=\frac1{\sqrt2}
					\left[\phi_1(x)e^{-iE_1t/\hbar}
					+\phi_2(x)e^{-iE_2t/\hbar}\right]}.
				\]

				\item 利用
				\[
					\langle n|x|n\rangle=\frac a2,\qquad
					\langle1|x|2\rangle=-\frac{16a}{9\pi^2},
				\]
				以及 $\langle1|p|2\rangle$ 为纯虚数，得到
				\[
					\boxed{\langle x\rangle_{t=0}
					=\frac a2-\frac{16a}{9\pi^2}},\qquad
					\boxed{\langle p\rangle_{t=0}=0}.
				\]

				\item 能量可取 $E_1$ 和 $E_2=4E_1$，概率均为 $1/2$，且与时间无关。因此
				\[
					\boxed{\langle H\rangle=\frac{E_1+E_2}{2}
					=\frac{5\pi^2\hbar^2}{4ma^2}}.
				\]
			\end{enumerate}
		}
	
\item 
考虑 $\mid lm>$ 态矢量空间中 $l=1$ 的子空间。在 $\left(\hat{l}^{2}, \hat{l}_{z}\right)$ 表象中角动量分量算符的矩阵表示为：
$$
\mathrm{L}_{\mathrm{x}}=\frac{\hbar}{\sqrt{2}}\left[\begin{array}{lll}
	0 & 1 & 0 \\
	1 & 0 & 1 \\
	0 & 1 & 0
\end{array}\right], \mathrm{L}_{\mathrm{y}}=\frac{\hbar}{\sqrt{2}}\left[\begin{array}{ccc}
	0 & -\mathrm{i} & 0 \\
	\mathrm{i} & 0 & -\mathrm{i} \\
	0 & \mathrm{i} & 0
\end{array}\right], \mathrm{L}_{\mathrm{z}}=\hbar\left[\begin{array}{ccc}
	1 & 0 & 0 \\
	0 & 0 & 0 \\
	0 & 0 & -1
\end{array}\right]
$$
\begin{enumerate}
	%\renewcommand{\labelenumi}{\arabic{enumi}.}
	% A(\Alph) a(\alph) I(\Roman) i(\roman) 1(\arabic)
	%设定全局标号series=example	%引用全局变量resume=example
	%[topsep=-0.3em,parsep=-0.3em,itemsep=-0.3em,partopsep=-0.3em]
	%可使用leftmargin调整列表环境左边的空白长度 [leftmargin=0em]
	\item
	求 $\mathrm{L}_{\mathrm{x}}$ 矩阵的本征值和本征矢;
	\item 
	求联系 $\left(\hat{l}^{2}, \hat{l}_{z}\right)$表象和$\left(\hat{l}^{2}, \hat{l}_{x}\right)$ 表象的幺正变换矩阵 $\mathrm{S}$;
	\item 
	利用幺正变换矩阵 $\mathrm{S}$, 求出 $\left(\hat{l}^{2}, \hat{l}_{x}\right)$ 表象中的矩阵表示 $\hat{l}_{x}, \hat{l}_{y}, \hat{l}_{z}$ 。
	
	
	
\end{enumerate}
\banswer{
	在 $(L^2,L_z)$ 表象中，基底均按磁量子数
	$+1,0,-1$ 排列。
	\begin{enumerate}
		\item 写 $L_x=\hbar M/\sqrt2$，特征方程为
		$\det(M-\eta I)=-\eta(\eta^2-2)=0$，所以本征值为 $+\hbar,0,-\hbar$。
		对 $\eta=\pm\sqrt2$，方程给出 $v_1=v_3$、$v_2=\eta v_1$；
		对 $\eta=0$，给出 $v_2=0$、$v_1=-v_3$。归一化得
		\[
			v_+=\begin{pmatrix}\frac12\\[2pt]\frac1{\sqrt2}\\[2pt]\frac12\end{pmatrix},
			\quad
			v_0=\begin{pmatrix}\frac1{\sqrt2}\\[2pt]0\\[2pt]-\frac1{\sqrt2}\end{pmatrix},
			\quad
			v_-=\begin{pmatrix}\frac12\\[2pt]-\frac1{\sqrt2}\\[2pt]\frac12\end{pmatrix}.
		\]

		\item 以这三个本征矢为列，得到从 $L_x$ 表象到 $L_z$ 表象的幺正矩阵
		\[
			\boxed{S=\begin{pmatrix}
			\frac12&\frac1{\sqrt2}&\frac12\\
			\frac1{\sqrt2}&0&-\frac1{\sqrt2}\\
			\frac12&-\frac1{\sqrt2}&\frac12
			\end{pmatrix}},\qquad S^\dagger S=I.
		\]

		\item 在 $L_x$ 表象中 $O_x=S^\dagger O_zS$，从而
		\[
			\boxed{L_x=\hbar\begin{pmatrix}1&0&0\\0&0&0\\0&0&-1\end{pmatrix}},
		\]
		\[
			\boxed{L_y=\frac{\hbar}{\sqrt2}
			\begin{pmatrix}0&i&0\\-i&0&i\\0&-i&0\end{pmatrix}},
			\qquad
			\boxed{L_z=\frac{\hbar}{\sqrt2}
			\begin{pmatrix}0&1&0\\1&0&1\\0&1&0\end{pmatrix}}.
		\]
	\end{enumerate}
}

\item 
一个磁矩为 $\vec{\mu}=\mu_{0} \vec{\sigma}$ 的自旋 $1 / 2$ 粒子在 $\mathrm{t}=0$ 时处于沿 $\mathrm{z}$ 正方向的均匀磁场 $\overrightarrow{\mathrm{B}}_{0}$
中, 其中$ \overrightarrow{\sigma} $为 Pauli 算符。当 $\mathrm{t}>0$ 时再加上一个旋转磁场 $\vec{B}_{1}(\mathrm{t})$, 其方向和 $z$ 轴垂	
直: $\overrightarrow{\mathrm{B}}_{1}(t)=\mathrm{B}_{1} \cos \left(2 \omega_{0} t \right) \overrightarrow{\mathrm{e}}_{x}-\mathrm{B}_{1} \sin \left(2 \omega_{0} t\right) \overrightarrow{\mathrm{e}}_{y}$, 其中 $\omega_{0}=\frac{\mu_{0} B_{0}}{\hbar}$ 。
\begin{enumerate}
	%\renewcommand{\labelenumi}{\arabic{enumi}.}
	% A(\Alph) a(\alph) I(\Roman) i(\roman) 1(\arabic)
	%设定全局标号series=example	%引用全局变量resume=example
	%[topsep=-0.3em,parsep=-0.3em,itemsep=-0.3em,partopsep=-0.3em]
	%可使用leftmargin调整列表环境左边的空白长度 [leftmargin=0em]
	\item
	求 $t>0$ 时的哈密顿算符 $\mathrm{H}=-\vec{\mu} \cdot\left[\overrightarrow{\mathrm{B}}_{0}+\overrightarrow{\mathrm{B}}_{1}({t})\right]$ 的本征值。 
	\item 
	若当 $t=0$ 时刻测得其自旋向下，求 $t$ 时刻自旋向上的概率。
	
	
	
	
\end{enumerate}
\banswer{
	令 $\omega_1=\mu_0B_1/\hbar$。在 $S_z$ 表象中
	\[
		H(t)=-\hbar\begin{pmatrix}
		\omega_0&\omega_1e^{i2\omega_0t}\\
		\omega_1e^{-i2\omega_0t}&-\omega_0
		\end{pmatrix}.
	\]
	\begin{enumerate}
		\item 因 $(\vec B_0+\vec B_1)^2=B_0^2+B_1^2$，瞬时本征值为
		\[
			\boxed{E_\pm=\pm\mu_0\sqrt{B_0^2+B_1^2}
			=\pm\hbar\sqrt{\omega_0^2+\omega_1^2}}.
		\]

		\item 旋转场的角频率 $2\omega_0$ 恰等于静磁场造成的两能级劈裂频率。令 $R(t)=e^{i\omega_0t\sigma_z}$，写 $|\psi(t)\rangle=R(t)|\psi_{\rm rot}(t)\rangle$。
		由 $R\sigma_xR^\dagger$ 的非对角元为 $e^{\pm2i\omega_0t}$，
		\[
		 H_{\rm rot}=R^\dagger HR-i\hbar R^\dagger\dot R
		 =-\hbar\omega_0\sigma_z-\hbar\omega_1\sigma_x+\hbar\omega_0\sigma_z
		 =-\hbar\omega_1\sigma_x.
		\]
		由 $\sigma_x^2=I$，演化算符为
		$e^{i\omega_1t\sigma_x}=I\cos(\omega_1t)+i\sigma_x\sin(\omega_1t)$。
		从 $|\downarrow\rangle$ 出发，
		\[
			|\psi_{\rm rot}(t)\rangle
			=\cos(\omega_1t)|\downarrow\rangle
			+i\sin(\omega_1t)|\uparrow\rangle.
		\]
		转回实验室表象只改变两个分量的相位，故
		\[
			\boxed{P_{\downarrow\to\uparrow}(t)=\sin^2(\omega_1t)
			=\sin^2\left(\frac{\mu_0B_1t}{\hbar}\right)}.
		\]
	\end{enumerate}
}


\newpage
\item 
一个质量为 $m$ 带电为$ -e  $的粒子在 $X-Y$ 平面内，做半径为 $r$ 的圆周运动。
\begin{enumerate}
	%\renewcommand{\labelenumi}{\arabic{enumi}.}
	% A(\Alph) a(\alph) I(\Roman) i(\roman) 1(\arabic)
	%设定全局标号series=example	%引用全局变量resume=example
	%[topsep=-0.3em,parsep=-0.3em,itemsep=-0.3em,partopsep=-0.3em]
	%可使用leftmargin调整列表环境左边的空白长度 [leftmargin=0em]
	\item
	求该粒子做圆周运动的能量本征值和本征函数。
	\item 
	 若在外场中添加一个微扰电场 $\overrightarrow{\mathscr{E}}=-\mathscr{E} \overrightarrow{\mathrm{e}}_{x}$，运用微扰理论, 求第一、第二激 发能量的二级近似。
	
	
	
\end{enumerate}

\banswer{
	以圆周角坐标 $\varphi$ 描述粒子，并记
	$\varepsilon=\hbar^2/(2mr^2)$。
	\begin{enumerate}
		\item
		\[
			H_0=-\frac{\hbar^2}{2mr^2}\frac{d^2}{d\varphi^2},
			\qquad
			\boxed{\psi_n(\varphi)=\frac{e^{in\varphi}}{\sqrt{2\pi}}},
			\qquad
			\boxed{E_n^{(0)}=\varepsilon n^2},
		\]
		其中 $n=0,\pm1,\pm2,\ldots$。除 $n=0$ 外，$\pm n$ 二重简并。

		\item 对 $\vec{\mathscr E}=-\mathscr E\,\vec e_x$ 可取
		$\Phi=\mathscr E r\cos\varphi$，故电荷 $-e$ 的微扰为
		$H'=-e\mathscr E r\cos\varphi$。记 $C=e\mathscr E r$，则
		\[
			\langle n'|H'|n\rangle=-\frac C2
			(\delta_{n',n+1}+\delta_{n',n-1}).
		\]
		在 $n=\pm1$ 和 $n=\pm2$ 各子空间中，一级矩阵均为零。
		二级有效矩阵的定义为
		\[
		 W^{(2)}_{ab}=\sum_{k\notin\mathcal D}
		 \frac{H'_{ak}H'_{kb}}{E^{(0)}-E_k^{(0)}}.
		\]
		对 $n=1$，对角元只有 $k=0,2$ 两项：
		$C^2[1/\varepsilon-1/(3\varepsilon)]/4=C^2/(6\varepsilon)$。
		$1$ 与 $-1$ 之间只有共同中间态 $k=0$，故非对角元为 $C^2/(4\varepsilon)$。
		因此第一激发子空间的二级有效矩阵为
		\[
			\frac{C^2}{\varepsilon}
			\begin{pmatrix}\frac16&\frac14\\[2pt]\frac14&\frac16\end{pmatrix}.
		\]
		所以由 $\cos\varphi$、$\sin\varphi$ 描写的两个实组合分别有
		\[
			\boxed{E_{1,\cos}=\varepsilon+\frac{5C^2}{12\varepsilon}},
			\qquad
			\boxed{E_{1,\sin}=\varepsilon-\frac{C^2}{12\varepsilon}}.
		\]
		对 $n=2$，对角元来自 $k=1,3$：
		$C^2[1/(3\varepsilon)-1/(5\varepsilon)]/4=C^2/(30\varepsilon)$。
		$2$ 与 $-2$ 不能通过一次中间跃迁相连，非对角元为零，所以到二级不分裂，二者均为
		\[
			\boxed{E_{2,\pm}=4\varepsilon+\frac{C^2}{30\varepsilon}}.
		\]
	\end{enumerate}
}


\item 
质量为 $m$ 的粒子在长度为 $a$ 的一维无限深对称方势阱 $(|x|<a / 2)$ 中处于
基态。设 $t=0$ 时刻突然添加一高度为 $c$, 宽度为 $b(<a)$ 的方形微扰势 $H^{\prime}$, 在 $t=T$ 时刻突然撤去该微扰势。
\begin{enumerate}
	%\renewcommand{\labelenumi}{\arabic{enumi}.}
	% A(\Alph) a(\alph) I(\Roman) i(\roman) 1(\arabic)
	%设定全局标号series=example	%引用全局变量resume=example
	%[topsep=-0.3em,parsep=-0.3em,itemsep=-0.3em,partopsep=-0.3em]
	%可使用leftmargin调整列表环境左边的空白长度 [leftmargin=0em]
	\item
	写出一维无限深对称方势阱哈密顿算符 $\mathrm{H}_{0}$ 的本征值和归一化本征函数。
	\item 
	写出 $\mathrm{t}>\mathrm{T}$ 时刻态矢量 $|\varphi(\mathrm{t})\rangle$ 的表达式。
	\item 
	试用一级微扰近似计算在 $t>\mathrm{T}$ 后体系跃迁到第一、第二、第三激发态的概率，并算出当微扰势的宽度$ b $取何值时对应跃迁概率取最大值。
	
	
	
\end{enumerate}

\banswer{
	假定题中的方形脉冲位于势阱中央，即
	$H'(x,t)=c$（$|x|<b/2,\,0<t<T$），其余处为零。
	\begin{enumerate}
		\item 令 $\epsilon=\pi^2\hbar^2/(2ma^2)$。本征值为
		\[
			E_n=n^2\epsilon,\qquad n=1,2,\ldots,
		\]
		归一化本征函数可取
		\[
			\phi_n(x)=\sqrt{\frac2a}
			\begin{cases}
			\cos(n\pi x/a),&n\ {\rm odd},\\
			\sin(n\pi x/a),&n\ {\rm even},
			\end{cases}
			\qquad |x|<\frac a2.
		\]

		\item 撤去势垒后的精确演化形式为
		\[
		 |\varphi(t)\rangle=e^{-iH_0(t-T)/\hbar}e^{-i(H_0+H')T/\hbar}|1\rangle.
		\]
		在相互作用表象展开到一级，得到
		\[
		 |\varphi(t)\rangle=\left(1-\frac{iH'_{11}T}{\hbar}\right)
		 e^{-iE_1t/\hbar}|1\rangle
		 +\sum_{n\ne1}c_n^{(1)}(T)e^{-iE_nt/\hbar}|n\rangle+O(c^2),
		\]
		其中
		\[
			c_n^{(1)}(T)=\frac{H'_{n1}}{\hbar\omega_{n1}}
			\left(1-e^{i\omega_{n1}T}\right),\qquad
			\omega_{n1}=\frac{E_n-E_1}{\hbar}.
		\]
		此处 $c_n^{(1)}$ 是相互作用表象系数，因此必须乘 $e^{-iE_nt/\hbar}$。
		若改用 $e^{-iE_n(t-T)/\hbar}$，系数就必须先乘 $e^{-iE_nT/\hbar}$。
		其中 $H'_{11}=c[b/a+\sin(\pi b/a)/\pi]$ 给出基态的一级附加相位。

		\item 中心微扰为偶函数，故从偶宇称基态只能跃迁到偶宇称态。因此第一、第三激发态（$n=2,4$）概率为
		\[
			\boxed{P_2=P_4=0}.
		\]
		对第二激发态 $n=3$，
		\[
			H'_{31}=\frac c\pi\left[
			\sin\frac{\pi b}{a}+\frac12\sin\frac{2\pi b}{a}\right],
		\]
		所以
		\[
			\boxed{P_3=
			\frac{4|H'_{31}|^2}{\hbar^2\omega_{31}^2}
			\sin^2\frac{\omega_{31}T}{2}},
			\qquad
			\omega_{31}=\frac{4\pi^2\hbar}{ma^2}.
		\]
		令 $u=\pi b/a$，空间因子为
		$\sin u(1+\cos u)$；其导数为
		$\cos u+\cos2u=2\cos^2u+\cos u-1$。
		内点极值满足 $\cos u=1/2$，而两端空间因子均为零，故
		\[
			\boxed{b=\frac a3}
		\]
		时第二激发态跃迁概率最大。两个恒为零的奇宇称跃迁不存在唯一的最优 $b$。
	\end{enumerate}
}


\end{enumerate}
